\begin{proposition}[2.1.4]
\label{IV.2.1.4-fr}
Sean $f:X\to Y$, $g:Y'\to Y$ dos morfismos de preesquemas, $\mathscr F$
un $\mathcal O_X$-módulo cuasicoherente; pongamos $X'=X\times_Y Y'$,
$f'=f_{(Y')}:X'\to Y'$, $\mathscr F'=\mathscr F\otimes_{\mathcal O_Y}\mathcal O_{Y'}$
y sea $g'$ la proyección canónica $X'\to X$. Sea $x'$ un punto de
$X'$, $x=g'(x')$, $y'=f'(x')$, $y=g(y')=f(x)$. Si $\mathscr F$ es $f$-plano
en el punto $x$, $\mathscr F'$ es $f'$-plano en el punto $x'$; en particular, si
$\mathscr F$ es $f$-plano, $\mathscr F'$ es $f'$-plano; si $f$ es plano,
$f'$ es plano.
\end{proposition}

Basta demostrar la primera afirmación; aplicando tres veces (\textbf I, 3.6.5),
así como (\textbf I, 2.4.4), es posible reducirse al caso en que
$Y=\operatorname{Spec}(\mathcal O_y)$, $X=\operatorname{Spec}(\mathcal O_x)$,
$Y'=\operatorname{Spec}(\mathcal O_{y'})$, $\mathscr F=\widetilde M$, donde
$M=\mathscr F_x$; la hipótesis y (2.1.2) implican entonces que $\mathscr F$
es $f$-plano, es decir, la cuestión se reduce a demostrar un caso particular de la
segunda afirmación, y esta última se deduce inmediatamente de (2.1.3).

\begin{proposition}[2.1.5]
\label{IV.2.1.5-fr}
Consideremos un diagrama conmutativo de morfismos de preesquemas
\[
\xymatrix{
  X \ar[d]_f & X' \ar[l]_{g'} \ar[d]^{f'} \\
  Y & Y' \ar[l]^g \ar[d]^h \\
  & Z
}
\]
donde $X'=X\times_Y Y'$ y $f'=f_{(Y')}$. Sea $x'$ un punto de $X'$, y
pongamos $x=g'(x')$, $y'=f'(x')$, $y=f(x)=g(y')$, $z=h(y')$. Sean
$\mathscr F$ un $\mathcal O_X$-módulo cuasicoherente $f$-plano en el punto $x$
(\textup{resp.} $f$-plano), y $\mathscr G'$ un $\mathcal O_{Y'}$-módulo cuasicoherente
$h$-plano en el punto $y'$ (\textup{resp.} $h$-plano); entonces
$\mathscr F\otimes_{\mathcal O_{Y'}}\mathscr G'$ es un $\mathcal O_{X'}$-módulo
cuasicoherente $(h\circ f')$-plano en el punto $x'$ (\textup{resp.} $(h\circ f')$-plano).
\end{proposition}

Como en (2.1.4), la cuestión se reduce al caso en que
$X=\operatorname{Spec}(\mathcal O_x)$, $Y=\operatorname{Spec}(\mathcal O_y)$,
\oldpage[IV-2]{7}
$Y'=\operatorname{Spec}(\mathcal O_{y'})$ y
$Z=\operatorname{Spec}(\mathcal O_z)$, y basta entonces demostrar que
$\mathscr F\otimes_{\mathcal O_{Y'}}\mathscr G'$ es $(h\circ f')$-plano. Teniendo en cuenta
(2.1.2), la proposición se deduce de Bourbaki, \emph{Alg. comm.},
chap.~I\textsuperscript{er}, \S~2, n\textsuperscript{o}~7, prop.~8.

\begin{corollary}[2.1.6]
\label{IV.2.1.6-fr}
Sean $f:X\to Y$, $g:Y\to Z$ dos morfismos de preesquemas, $\mathscr F$ un
$\mathcal O_X$-módulo. Si $\mathscr F$ es $f$-plano en el punto $x\in X$ y si
$g$ es plano en el punto $f(x)$, $\mathscr F$ es $(g\circ f)$-plano en el punto $x$. En
particular, si $f$ y $g$ son morfismos planos, lo mismo vale para
$g\circ f$.
\end{corollary}

Es el caso particular de (2.1.5) con $Y'=Y$, $\mathscr G'=\mathcal O_{Y'}$.

\begin{corollary}[2.1.7]
\label{IV.2.1.7-fr}
Si $f:X\to X'$, $g:Y\to Y'$ son dos $S$-morfismos planos,
$f\times_S g:X\times_S Y\to X'\times_S Y'$ es un morfismo plano.
\end{corollary}

Esto se deduce de (2.1.4) y (2.1.6) (cf. \textbf I, 3.5.1).
